\subsection{Standard Distribution Families and Approximations}
\begin{definition} \label{def:core-discrete-distributions} \hlt{Core Discrete Distributions}\\
Common discrete laws include:
\begin{enumerate}[label=(\roman*)]
\setlength{\itemsep}{0pt}
\item \hlt{Point mass}: $X\sim\delta_a$ for $a\in\R$ if $\mathbb{P}(X=a)=1$.
\item \hlt{Discrete uniform}: $X\sim\operatorname{Unif}\{1,\ldots,k\}$ for $k\in\N$ if $\mathbb{P}(X=x)=\frac{1}{k}$ for $x=1,\ldots,k$.
\item \hlt{Bernoulli}: $X\sim\operatorname{Bernoulli}(p)$ for $p\in[0,1]$ if $\mathbb{P}(X=1)=p$ and $\mathbb{P}(X=0)=1-p$.\\
Equivalently $\mathbb{P}(X=x)=p^x(1-p)^{1-x}$ for $x\in\{0,1\}$.
\item \hlt{Binomial}: $X\sim\operatorname{Binomial}(n,p)$ for $n\in\N$ and $p\in[0,1]$ if $\mathbb{P}(X=x)=\binom{n}{x}p^x(1-p)^{n-x}$ for $x=0,\ldots,n$.
\item \hlt{Geometric}: $X\sim\operatorname{Geom}(p)$, $p\in(0,1]$, if $\mathbb{P}(X=k)=p(1-p)^{k-1}$ for $k\geq1$; $X$ counts trials to first success.
\item \hlt{Negative binomial}: $X\sim\operatorname{NegBin}(r,p)$, $r\in\N$, $p\in(0,1]$, if $\mathbb{P}(X=k)=\binom{k-1}{r-1}p^r(1-p)^{k-r}$ for $k\geq r$;\\
$X$ counts trials to the $r$th success.
\item \hlt{Hypergeometric}: $X\sim\operatorname{Hypergeom}(n,r,m)$ if $\mathbb{P}(X=k)=\frac{\binom{r}{k}\binom{n-r}{m-k}}{\binom{n}{m}}$\\
for $\max\{0,m-n+r\}\leq k\leq\min\{r,m\}$.
\item \hlt{Poisson}: $X\sim\operatorname{Poisson}(\lambda)$ for $\lambda\geq0$ if $\mathbb{P}(X=x)=\frac{e^{-\lambda}\lambda^x}{x!}$ for $x=0,1,2,\ldots$.
\end{enumerate}
\end{definition}

\begin{definition} \label{def:core-continuous-distributions} \hlt{Core Continuous Distributions}\\
Common continuous laws include:
\begin{enumerate}[label=(\roman*)]
\setlength{\itemsep}{0pt}
\item \hlt{Uniform}: $X\sim\operatorname{Uniform}(a,b)$ for $a<b$ if $f_X(x)=(b-a)^{-1}\mathbf{1}_{[a,b]}(x)$.
\item \hlt{Normal}: $X\sim N(\mu,\sigma^2)$ for $\mu\in\R$, $\sigma>0$, if $f_X(x)=(\sigma\sqrt{2\pi})^{-1}\exp(-\frac{(x-\mu)^2}{2\sigma^2})$ for $x\in\R$.\\
$N(0,1)$ has density $\phi$ and CDF $\Phi$.
\item \hlt{Exponential}: $X\sim\operatorname{Exp}(\lambda)$ for $\lambda>0$ if $f_X(x)=\lambda e^{-\lambda x}\mathbf{1}_{[0,\infty)}(x)$.
\item \hlt{Gamma}: $X\sim\operatorname{Gamma}(\alpha,\lambda)$ for $\alpha,\lambda>0$ if $f_X(x)=\lambda^\alpha\Gamma(\alpha)^{-1}x^{\alpha-1}e^{-\lambda x}\mathbf{1}_{[0,\infty)}(x)$; $\operatorname{Exp}(\lambda)=\operatorname{Gamma}(1,\lambda)$.
\item \hlt{Beta}: $X\sim\operatorname{Beta}(\alpha,\beta)$ for $\alpha,\beta>0$ if $f_X(x)=\Gamma(\alpha+\beta)(\Gamma(\alpha)\Gamma(\beta))^{-1}x^{\alpha-1}(1-x)^{\beta-1}\mathbf{1}_{(0,1)}(x)$.
\item \hlt{Cauchy}: $X\sim\operatorname{Cauchy}(0,1)$ if $f_X(x)=(\pi(1+x^2))^{-1}$ for $x\in\R$.
\item \hlt{$\chi^2$}: $X\sim\chi_p^2$ for $p>0$ if $f_X(x)=(\Gamma(\frac{p}{2})2^{\frac{p}{2}})^{-1}x^{\frac{p}{2}-1}e^{-\frac{x}{2}}\mathbf{1}_{(0,\infty)}(x)$.
\end{enumerate}
\end{definition}

\begin{definition} \label{def:normal-right-tail-quantile} \hlt{Normal Right-Tail Quantile}\\
For $0<\alpha<1$ and $Z\sim N(0,1)$, $z(\alpha)\defeq\Phi^{-1}(1-\alpha)$, so $\mathbb{P}(Z>z(\alpha))=\alpha$.
\end{definition}

\begin{method} \label{meth:discrete-distribution-model-recipes} \hlt{Discrete Distribution Model Recipes}
\begin{enumerate}[label=(\roman*)]
\setlength{\itemsep}{0pt}
\item \hlt{Indicators}: for $A\in\mathcal{F}$, $\mathbf{1}_A\sim\operatorname{Bernoulli}(\mathbb{P}(A))$.
\item \hlt{Binomial counts}: if $B_1,\ldots,B_n$ are independent $\operatorname{Bernoulli}(p)$ variables, then $\sum\limits_{i=1}^{n}B_i\sim\operatorname{Binomial}(n,p)$.
\item \hlt{Geometric waiting}: the probability of first success on trial $k$ is $(1-p)^{k-1}p$.
\item \hlt{Negative-binomial waiting}: for iid $B_i\sim\operatorname{Bernoulli}(p)$ and $T_r=\min\{k:\sum\limits_{i=1}^{k}B_i=r\}$,\\
$\{T_r=k\}=\{B_k=1,\sum\limits_{i=1}^{k-1}B_i=r-1\}$, so $\mathbb{P}(T_r=k)=\binom{k-1}{r-1}p^r(1-p)^{k-r}$.
\item \hlt{Hypergeometric sampling}: if $M\subseteq[n]$, $\abs{M}=r$, and $S$ is a uniform $m$-subset,\\
then $X=\abs{S\cap M}$ and $\mathbb{P}(X=k)=\frac{\binom{r}{k}\binom{n-r}{m-k}}{\binom{n}{m}}$ for $\max\{0,m-n+r\}\leq k\leq\min\{r,m\}$.
\end{enumerate}
\hlt{Key Insight}: Name the event; PMFs come from independence or subset counts.
\end{method}

\begin{proposition} \label{prop:poisson-limit-of-binomial} \hlt{Poisson Limit of Rare Binomial Counts}\\
Let $\lambda>0$, $k\in\N_0$, and $X_n\sim\operatorname{Binomial}(n,\frac{\lambda}{n})$ for all sufficiently large $n$. Then $\lim\limits_{n\to\infty}\mathbb{P}(X_n=k)=\frac{e^{-\lambda}\lambda^k}{k!}$.
\begin{proof}
For large $n$, $\mathbb{P}(X_n=k)=\binom{n}{k}(\frac{\lambda}{n})^k(1-\frac{\lambda}{n})^{n-k}=\frac{\lambda^k}{k!}\prod\limits_{j=0}^{k-1}(1-\frac{j}{n})(1-\frac{\lambda}{n})^n(1-\frac{\lambda}{n})^{-k}$.\\
Since $\prod\limits_{j=0}^{k-1}(1-\frac{j}{n})\to1$, $(1-\frac{\lambda}{n})^n\to e^{-\lambda}$, and $(1-\frac{\lambda}{n})^{-k}\to1$, the limit is $\frac{e^{-\lambda}\lambda^k}{k!}$.\\
\hlt{Structural Commentary}: Rare binomial masses split into fixed-$k$ choices and an exponential no-event factor.\\
\hlt{Key Insight}: Large $n$, small $p$, fixed $np$ turns binomial masses into Poisson masses.
\end{proof}
\end{proposition}

\begin{method} \label{meth:recursive-binomial-poisson-masses} \hlt{Recursive Binomial and Poisson Mass Computation}\\
Let $q\defeq1-p$.
\begin{enumerate}[label=(\roman*)]
\setlength{\itemsep}{0pt}
\item If $X\sim\operatorname{Binomial}(n,p)$, $p_k\defeq\mathbb{P}(X=k)$, and $q>0$, then $p_0=q^n$ and $\frac{p_k}{p_{k-1}}=\frac{(n-k+1)p}{kq}$ for $1\leq k\leq n$.
\item If $X\sim\operatorname{Poisson}(\lambda)$ and $p_k\defeq\mathbb{P}(X=k)$, then $p_0=e^{-\lambda}$ and $\frac{p_k}{p_{k-1}}=\frac{\lambda}{k}$ for $k\geq1$.
\end{enumerate}
\hlt{Key Insight}: Successive mass ratios avoid repeatedly evaluating large factorials.
\end{method}

\begin{example} \label{ex:majority-decoder-binomial-tail} \hlt{Majority Decoding as a Binomial Tail}\\
Let $n=2m+1$ and let $C_i\sim\operatorname{Bernoulli}(p)$ be independent corruption indicators.\\
Then $X=\sum\limits_{i=1}^{n}C_i\sim\operatorname{Binomial}(n,p)$ and the majority decoder is correct iff $X\leq m$.\\
For $n=5$, $p=0.1$, $\mathbb{P}(X\leq2)=0.9^5+5(0.1)(0.9)^4+10(0.1)^2(0.9)^3\approx0.9914$.\\
\hlt{Key Insight}: Repetition coding converts reliability into a lower binomial tail.
\end{example}

\begin{example} \label{ex:poisson-approximation-double-sixes} \hlt{Poisson Approximation for Rare Repeated Events}\\
Let $I_j=\mathbf{1}_{\{\text{double six on roll }j\}}$ for $1\leq j\leq100$.\\
Then $I_j\sim\operatorname{Bernoulli}(\frac{1}{36})$ iid, $X=\sum\limits_{j=1}^{100}I_j\sim\operatorname{Binomial}(100,\frac{1}{36})$, and with $\lambda=\frac{100}{36}=\frac{25}{9}$, $\mathbb{P}(X=k)\approx \frac{e^{-\frac{25}{9}}(\frac{25}{9})^k}{k!}$.\\
\hlt{Key Insight}: Rare repeated event counts are often well approximated by a Poisson law with $\lambda=np$.
\end{example}

\begin{proposition} \label{prop:exponential-memoryless-property} \hlt{Exponential Memorylessness}\\
If $T\sim\operatorname{Exp}(\lambda)$ with $\lambda>0$, then $\mathbb{P}(T>s+t\mid T>s)=\mathbb{P}(T>t)$ for all $s,t\geq0$.
\begin{proof}
For $u\geq0$, $\mathbb{P}(T>u)=e^{-\lambda u}$. Hence $\mathbb{P}(T>s+t\mid T>s)=\frac{\mathbb{P}(T>s+t)}{\mathbb{P}(T>s)}=\frac{e^{-\lambda(s+t)}}{e^{-\lambda s}}=e^{-\lambda t}=\mathbb{P}(T>t)$.\\
\hlt{Structural Commentary}: Conditioning on survival to time $s$ leaves the residual tail unchanged.\\
\hlt{Key Insight}: Exponential tails convert elapsed time into a cancellable factor.
\end{proof}
\end{proposition}

\begin{remark} \label{rem:poisson-process-distribution-link} \hlt{Poisson Counts and Exponential Gaps}\\
In homogeneous Poisson process, counts satisfy $N(t)\sim\operatorname{Poisson}(\lambda t)$; interarrival gaps satisfy $S_i\sim\operatorname{Exp}(\lambda)$ iid.
\end{remark}

\begin{proposition} \label{prop:normal-standardization-and-sums} \hlt{Standardizing Normal Variables}\\
For $\sigma,\sigma_i>0$:
\begin{enumerate}[label=(\roman*)]
\setlength{\itemsep}{0pt}
\item \hlt{Standardization}: if $X\sim N(\mu,\sigma^2)$, then $\frac{X-\mu}{\sigma}\sim N(0,1)$.
\item \hlt{Affine reconstruction}: if $Z\sim N(0,1)$, then $\mu+\sigma Z\sim N(\mu,\sigma^2)$.
\item \hlt{Independent sums}: if $X_i\sim N(\mu_i,\sigma_i^2)$ are independent, then $\sum\limits_{i=1}^{n}X_i\sim N(\sum\limits_{i=1}^{n}\mu_i,\sum\limits_{i=1}^{n}\sigma_i^2)$.
\item \hlt{Interval probabilities}: if $X\sim N(\mu,\sigma^2)$ and $a<b$, then $\mathbb{P}(a<X<b)=\Phi(\frac{b-\mu}{\sigma})-\Phi(\frac{a-\mu}{\sigma})$.
\end{enumerate}
\begin{proof}
\begin{enumerate}[label=(\roman*)]
\setlength{\itemsep}{0pt}
\item Let $Y=\frac{X-\mu}{\sigma}$. By change of variables, $f_Y(z)=\sigma f_X(\mu+\sigma z)=(2\pi)^{-\frac{1}{2}}e^{-\frac{z^2}{2}}$, so $Y\sim N(0,1)$.
\item Let $Y=\mu+\sigma Z$. Then $f_Y(y)=\sigma^{-1}\phi(\frac{y-\mu}{\sigma})=(\sigma\sqrt{2\pi})^{-1}\exp(-\frac{(y-\mu)^2}{2\sigma^2})$, hence $Y\sim N(\mu,\sigma^2)$.
\item It suffices by induction to treat $n=2$.\\
Let $X_i\sim N(\mu_i,v_i)$ be independent, $S=X_1+X_2$, $v=v_1+v_2$, $w=\frac{v_1v_2}{v}$, and $a_s=\frac{v_2\mu_1+v_1(s-\mu_2)}{v}$.\\
Then $f_S(s)=\int_\R f_1(u)f_2(s-u)\,du$, and completing the square gives\\
$f_1(u)f_2(s-u)=(2\pi v)^{-\frac{1}{2}}e^{-\frac{(s-\mu_1-\mu_2)^2}{2v}}(2\pi w)^{-\frac{1}{2}}e^{-\frac{(u-a_s)^2}{2w}}$.\\
The second factor integrates to $1$, so $S\sim N(\mu_1+\mu_2,v_1+v_2)$; induction gives (iii).
\item Let $Z=\frac{X-\mu}{\sigma}$. Since $\sigma>0$, $\{a<X<b\}=\{\frac{a-\mu}{\sigma}<Z<\frac{b-\mu}{\sigma}\}$.\\
Hence $\mathbb{P}(a<X<b)=\Phi(\frac{b-\mu}{\sigma})-\Phi(\frac{a-\mu}{\sigma})$.
\end{enumerate}
\hlt{Structural Commentary}: Normal laws are stable under affine maps and independent addition.\\
\hlt{Key Insight}: Standardize to $N(0,1)$, compute with $\Phi$, then transform back.
\end{proof}
\end{proposition}

\begin{example} \label{ex:normal-tail-and-quantile-standardization} \hlt{Normal Tail and Quantile Standardization}\\
Let $X\sim N(3,5)$ and $Z=\frac{X-3}{\sqrt5}\sim N(0,1)$. Then $\mathbb{P}(X>1)=\mathbb{P}(Z>-\frac{2}{\sqrt5})=1-\Phi(-\frac{2}{\sqrt5})\approx0.81$.\\
If $\mathbb{P}(X<q)=0.2$, then $\Phi(\frac{q-3}{\sqrt5})=0.2$, so $q=3+\sqrt5\,\Phi^{-1}(0.2)\approx3-0.8416\sqrt5\approx1.1181$.
\end{example}
